\chapter{What is a Generating Function?}

A sequence of numbers $(a_n)_{n \ge 0}$ is the natural output of a counting problem: how many binary strings of length $n$, how many trees with $n$ nodes, how many ways to parse a sentence of length $n$ according to a given grammar. Studying the sequence directly forces us to deal with each index separately, and recurrences --- while precise --- often resist closed-form extraction. The central idea of generating function theory is to encode the entire sequence as a single algebraic object,
\[
  A(z) \;=\; \sum_{n \ge 0} a_n z^n,
\]
and then let the algebra of formal power series do the bookkeeping. This is not mere notational convenience. Operations on sequences that would otherwise require index gymnastics --- shifting, convolving, differencing --- correspond to clean algebraic operations on $A(z)$. Recurrences, which are linear or polynomial relations among the $a_n$, become \emph{functional equations} for $A(z)$, and those equations can often be solved in closed form. The solution then exposes structure invisible in the raw sequence.

\begin{definition}
The ring of \emph{formal power series} over $\mathbb{Q}$, denoted $\mathbb{Q}[[z]]$, consists of all expressions $A(z) = \sum_{n \ge 0} a_n z^n$ with $a_n \in \mathbb{Q}$, with addition $(A+B)(z) = \sum_n (a_n + b_n) z^n$ and multiplication $(A \cdot B)(z) = \sum_n c_n z^n$ where $c_n = \sum_{k=0}^{n} a_k b_{n-k}$.
\end{definition}

At this stage $z$ is a \emph{formal indeterminate} --- it carries no numerical value, and questions of convergence are entirely irrelevant. The ring $\mathbb{Q}[[z]]$ is an integral domain, and any series $A(z)$ with $a_0 \ne 0$ is invertible: one constructs $A(z)^{-1}$ coefficient by coefficient from the equation $A \cdot A^{-1} = 1$. Two further operations are essential. The \emph{shift} by $z$ sends $A(z)$ to $z A(z) = \sum_{n \ge 1} a_{n-1} z^n$, so $[z^n](zA(z)) = a_{n-1}$. The \emph{formal derivative} is $A'(z) = \sum_{n \ge 1} n a_n z^{n-1}$, defined term by term with no analytic content; it satisfies the usual Leibniz rule $(\cdot\,\cdot)'=(\cdot)'\cdot + \cdot\,(\cdot)'$ by direct verification. The \emph{coefficient extraction operator} $[z^n] A(z) := a_n$ is linear and respects products via the convolution formula $[z^n](A \cdot B) = \sum_{k=0}^{n} a_k b_{n-k}$.

\section*{Binary strings}

Let $b_n$ denote the number of binary strings of length $n$. Every such string is a sequence of $n$ bits, each chosen from $\{0,1\}$, so $b_n = 2^n$. We can derive this from the recurrence and simultaneously introduce the method. Any binary string of length $n \ge 1$ is obtained by appending one of two symbols to a string of length $n-1$, giving $b_n = 2 b_{n-1}$ for $n \ge 1$, with $b_0 = 1$. Multiply both sides by $z^n$ and sum over $n \ge 1$:
\[
  \sum_{n \ge 1} b_n z^n \;=\; 2z \sum_{n \ge 1} b_{n-1} z^{n-1}.
\]
The left side is $B(z) - b_0 = B(z) - 1$, and the right side is $2z B(z)$. Thus
\[
  B(z) - 1 \;=\; 2z B(z), \qquad \text{so} \qquad B(z) \;=\; \frac{1}{1 - 2z}.
\]
This could equally be read as the geometric series $\sum_{n \ge 0} (2z)^n = \sum_{n \ge 0} 2^n z^n$, confirming $[z^n] B(z) = 2^n$. The functional equation route is more instructive: it works even when no closed-form formula for $b_n$ is obvious beforehand.

\section*{Fibonacci numbers}

The Fibonacci sequence is defined by $F_0 = 0$, $F_1 = 1$, and $F_n = F_{n-1} + F_{n-2}$ for $n \ge 2$. Setting $F(z) = \sum_{n \ge 0} F_n z^n$ and converting the recurrence: multiply $F_n = F_{n-1} + F_{n-2}$ by $z^n$ and sum over $n \ge 2$,
\[
  \sum_{n \ge 2} F_n z^n
  \;=\; z \sum_{n \ge 2} F_{n-1} z^{n-1}
  \;+\; z^2 \sum_{n \ge 2} F_{n-2} z^{n-2}.
\]
The left side is $F(z) - F_1 z - F_0 = F(z) - z$. The two sums on the right are $z$ times $\sum_{m \ge 1} F_m z^m = F(z) - F_0 = F(z)$ and $z^2$ times $\sum_{m \ge 0} F_m z^m = F(z)$, respectively. Collecting terms,
\[
  F(z) - z \;=\; z F(z) + z^2 F(z),
\]
which gives immediately
\[
  F(z) \;=\; \frac{z}{1 - z - z^2}.
\]
To extract $F_n$ explicitly, factor the denominator. The roots of $1 - z - z^2 = 0$ (equivalently, $z^2 + z - 1 = 0$) are $z = (-1 \pm \sqrt{5})/2$. In terms of the golden ratio $\varphi = (1 + \sqrt{5})/2$ and its conjugate $\hat{\varphi} = (1 - \sqrt{5})/2$, one checks that $1/\varphi$ and $1/\hat{\varphi}$ are the roots of $1 - z - z^2$ viewed as a polynomial in $z$, so
\[
  1 - z - z^2 \;=\; -\Bigl(z - \tfrac{1}{\varphi}\Bigr)\Bigl(z - \tfrac{1}{\hat\varphi}\Bigr).
\]
Partial fractions decompose $F(z)$ as
\[
  F(z) \;=\; \frac{1}{\sqrt{5}}\cdot\frac{1}{1 - \varphi z} \;-\; \frac{1}{\sqrt{5}}\cdot\frac{1}{1 - \hat\varphi z}.
\]
Expanding each geometric series and reading off the coefficient of $z^n$ gives Binet's formula,
\[
  F_n \;=\; \frac{\varphi^n - \hat\varphi^n}{\sqrt{5}}.
\]
The key observation is that this closed form was \emph{forced} by the location of the poles of $F(z)$. The dominant term as $n \to \infty$ comes from the pole closest to the origin, which is at $z = 1/\varphi$, yielding $F_n \sim \varphi^n / \sqrt{5}$. This is not a coincidence: singularities of generating functions control asymptotics, a theme developed fully in Chapters~3 and~4. The standard reference for the entire theory is \cite{FlajoletSedgewick2009}.

\section*{Catalan numbers}

The Catalan numbers $C_n$ count full binary trees (every internal node has exactly two children) with $n$ internal nodes; equivalently, they count triangulations of a convex $(n+2)$-gon, Dyck paths of length $2n$, and well-formed parenthesizations of length $n$, among many other combinatorial families. The sequence begins $1, 1, 2, 5, 14, 42, \ldots$ and satisfies the recurrence
\[
  C_n \;=\; \sum_{k=0}^{n-1} C_k\, C_{n-1-k}, \qquad C_0 = 1.
\]
This recurrence reflects a structural decomposition: a binary tree with $n \ge 1$ internal nodes consists of a root together with a left subtree on $k$ internal nodes and a right subtree on $n-1-k$ nodes, for some $0 \le k \le n-1$. The right-hand side is a convolution: $[z^{n-1}](C(z)^2) = \sum_{k=0}^{n-1} C_k C_{n-1-k}$. Consequently, multiplying by $z^{n-1}$ and summing over $n \ge 1$, the recurrence assembles into
\[
  C(z) \;=\; 1 + z\, C(z)^2.
\]
This is our first \emph{quadratic} functional equation. Solving for $C(z)$ using the quadratic formula,
\[
  C(z) \;=\; \frac{1 \pm \sqrt{1 - 4z}}{2z}.
\]
The sign must be chosen so that $C(z)$ is a valid formal power series with $C_0 = 1$. The plus sign gives $C(0) = 2/0$, which is undefined, so we take the minus sign:
\[
  C(z) \;=\; \frac{1 - \sqrt{1 - 4z}}{2z}.
\]
To verify and extract coefficients, expand $\sqrt{1-4z} = \sum_{n \ge 0} \binom{1/2}{n}(-4z)^n$ using the generalized binomial theorem (valid as a formal identity in $\mathbb{Q}[[z]]$). After simplification, one finds
\[
  [z^n] C(z) \;=\; \frac{1}{n+1}\binom{2n}{n},
\]
which is the standard formula for the $n$-th Catalan number. The appearance of $\sqrt{1-4z}$ is significant: $C(z)$ is not a rational function but an \emph{algebraic} one, satisfying a polynomial equation over $\mathbb{Q}(z)$. Algebraic generating functions have branch-point singularities rather than poles. The branch point of $C(z)$ at $z = 1/4$ is a square-root singularity, and from the local expansion $C(z) \sim 2 - 2\sqrt{1-4z}$ near $z = 1/4$, one can show (see Chapter~5) that
\[
  C_n \;\sim\; \frac{4^n}{n^{3/2}\sqrt{\pi}},
\]
a formula whose $n^{-3/2}$ subexponential factor is a direct signature of the square-root branch point. This pattern is universal: any combinatorial class whose generating function is algebraic with a square-root singularity as its dominant singularity will exhibit $n^{-3/2}$ subexponential decay in its coefficients. Unambiguous context-free grammars always produce algebraic generating functions, and this has concrete implications for the statistical structure of natural language sentences --- a connection we will revisit when discussing the relationship between pushdown automata and the token distributions produced by language models.

\section*{Ordinary versus exponential generating functions}

The generating functions above are \emph{ordinary generating functions} (OGFs): the sequence element $a_n$ is literally the coefficient of $z^n$. For \emph{unlabeled} combinatorial structures, where we only care about shape and not about which specific elements occupy which positions, OGFs are the natural choice. But many problems in combinatorics --- and in the analysis of language models --- involve \emph{labeled} structures, where the $n$ atoms in a structure are drawn from a distinguished set and the assignment of labels matters. The canonical example is permutations: the number of permutations of $n$ distinct elements is $n!$, and the product structure of permutations (concatenating two permutations of disjoint sets requires interleaving their labels, not merely concatenating them) is reflected not in ordinary convolution but in \emph{binomial convolution}: if $A$ and $B$ count labeled structures with $a_n$ and $b_n$ objects of size $n$, then the number of labeled structures formed by selecting a subset of $k$ labels for an $A$-structure and the complementary $n-k$ labels for a $B$-structure is $\sum_{k=0}^n \binom{n}{k} a_k b_{n-k}$, not $\sum_{k=0}^n a_k b_{n-k}$.

\begin{definition}
The \emph{exponential generating function} (EGF) of a sequence $(a_n)_{n \ge 0}$ is
\[
  \hat{A}(z) \;=\; \sum_{n \ge 0} a_n \frac{z^n}{n!}.
\]
\end{definition}

Binomial convolution of sequences corresponds to ordinary multiplication of EGFs: if $\hat{A}(z) = \sum a_n z^n/n!$ and $\hat{B}(z) = \sum b_n z^n/n!$, then $[z^n/n!](\hat A \cdot \hat B) = \sum_k \binom{n}{k} a_k b_{n-k}$, exactly the right combinatorial count. This is the algebraic reason that labeled structures prefer EGFs.

\begin{example}
Consider the sequence $a_n = n!$, counting permutations. Its OGF is $\sum_{n \ge 0} n! \, z^n$. This series has zero radius of convergence as an analytic function: the coefficients grow faster than any geometric sequence, so the series diverges for every $z \ne 0$. In the formal world of $\mathbb{Q}[[z]]$ it is a perfectly valid element, but its analytic utility is nil. Its EGF, on the other hand, is
\[
  \hat{S}(z) \;=\; \sum_{n \ge 0} n! \cdot \frac{z^n}{n!} \;=\; \sum_{n \ge 0} z^n \;=\; \frac{1}{1-z},
\]
which is a rational function with a simple pole at $z = 1$. Thus $[z^n/n!]\,\hat{S}(z) = 1$, and recovering the count of permutations of $[n]$ requires multiplying by $n!$: there are $1 \cdot n! = n!$ such permutations. The EGF encodes the sequence cleanly at the level of the combinatorial species.
\end{example}

The distinction between OGF and EGF is not merely aesthetic: using an OGF for a labeled structure or an EGF for an unlabeled one leads to incorrect functional equations. Part of the skill of analytic combinatorics is recognizing which generating function type matches the combinatorial product being modeled. The \emph{symbolic method}, developed in Chapter~2, makes this assignment systematic: the constructions $\SEQ$, $\SET$, $\CYC$, and $\MSET$ each translate directly into GF operations, and the labeled/unlabeled distinction determines whether the translation uses OGF or EGF conventions.

\section*{Road map}

This chapter has introduced generating functions purely as formal algebraic objects, and the examples have already shown their power: a recurrence becomes a functional equation, and the equation's solution encodes both exact counts and, implicitly, asymptotic behavior. The next chapter treats $z$ as a complex variable and develops the analytic theory that lets us extract asymptotics from singularities --- the principle that the growth rate of $a_n$ is governed by the location of the nearest singularity of $A(z)$ in $\mathbb{C}$, and the subexponential factor is governed by the type of that singularity. Chapter~3 introduces the symbolic method, turning combinatorial grammar rules directly into GF equations without any manual index manipulation. Chapters~4 and~5 then apply singularity analysis to rational and algebraic generating functions respectively, recovering the $\varphi^n$ growth of Fibonacci numbers and the $4^n n^{-3/2}$ growth of Catalan numbers from first principles. Throughout, we will draw connections to the statistical structure of sequences produced by large language models: such models define probability distributions over discrete sequences, and generating functions provide the right language for analyzing those distributions at scale.
